\noindent\textbf{Por qué se usa este procedimiento.} En cortante puro, $\sigma_x=\sigma_y=\sigma_z=0$. En la expresión de von Mises quedan tres veces los cuadrados de cortante; con una sola componente $\sigma_v^2=3\tau^2$. La raíz positiva da $\sqrt3|\tau|$.

\begin{enumerate}

\item Con $\sigma_x = \sigma_y = 0$, $\sigma_v^2 = 3\tau^2$.

\end{enumerate}

\noindent\textbf{Operaciones y resultados intermedios.}

\begin{enumerate}

\item Tensión de von Mises.

\[\fitmath{y_{1}=\sqrt{3}\cdot 100\approx 173{,}2051\,\mathrm{MPa}}\]

\end{enumerate}
